% FILE: 02_lineage_and_context.tex
% PROJECT: phd-thesis
% MANUFACTURED: 2026-06-01
% AUTHORITY: ledgers/PROJECT_INDEX.yaml (lineage header), checkpoints/PROCESS_INTELLIGENCE_ALIVE_001.md
% DO NOT CLAIM: See ledgers/DO_NOT_CLAIM_LEDGER.md

\section{Lineage and Context}
\label{sec:phd-thesis-lineage}

\subsection{2016 Language-Model Lesson}
\label{sec:phd-thesis-lineage-2016}

The thesis lineage recorded in the header of \texttt{ledgers/PROJECT\_INDEX.yaml} begins with
a 2016 language-model experiment. The finding of that experiment, stated in the lineage
record, is: \emph{local text imitation does not conserve consequence}. This is the
foundational negative result from which the entire research program derives.

The 2016 experiment is not documented in the 34 project manifests currently catalogued in the
foundry --- its source materials predate the \texttt{process-intelligence} repository. It is
cited here as lineage, not as a reproduced artifact: the experimental finding is the author's
thesis (AUTHOR THESIS), not a claim backed by a file in this repository. What the experiment
established was a boundary: language models trained to imitate local token sequences do not
learn the causal, temporal, and custodial structure of the processes those sequences describe.
A model that predicts the next word in an invoice approval log does not understand that
``Invoice Approval must follow Goods Receipt'' is a binding constraint --- it learns a
statistical correlation, not a process law.

This boundary --- the gap between local imitation and process consequence --- is the
motivating observation for the entire downstream research program. It explains why receipt
chains, conformance checking, and object-centric event logs are necessary rather than
optional: they exist to restore the consequence structure that prediction-only approaches
discard.

\subsection{Enterprise Process Gap}
\label{sec:phd-thesis-lineage-gap}

The enterprise process gap is the observable manifestation of the 2016 lesson at organizational
scale. It is the structural condition in which an enterprise's declared process model diverges
from its actual runtime process in ways that are invisible to the enterprise's own monitoring
and reporting systems. The gap surfaces only under formal process mining: discovery, conformance
checking, replay, and object-lifecycle analysis.

The adversarial research surface of the foundry documents four categories of this gap:
raw laundering refusal (emitting conformance claims without a named witness),
conformance theater (passing conformance gates on fabricated or incomplete logs),
metric fabrication (presenting fitness/precision scores without a verifiable model-log binding),
and variant explosion (declaring a single happy-path model while the mined log reveals
dozens of rework and exception variants). These findings are recorded in
\texttt{adversarial/RED\_TEAM\_FINDINGS\_001.md} and
\texttt{adversarial/RED\_TEAM\_FINDINGS\_002.md}, and are operationalized in the
adversarial scripts
\texttt{adversarial/simulate\_laundering\_spoofing.py} and
\texttt{adversarial/validate\_standards.py}.

The gap has a direct M\&A consequence: a buyer performing process diligence on an acquisition
target cannot rely on the target's internal conformance reports if those reports are not backed
by receipts. The M\&A claim taxonomy in \texttt{ma/BOARD\_CLAIM\_TAXONOMY.md} and the
slide-to-receipt map in \texttt{ma/SLIDE\_TO\_RECEIPT\_MAP.md} exist specifically to close this
gap at the board level --- converting process mining findings into board-admissible claims with
verifiable receipt chains.

\subsection{Relationship to the Chatman Equation}
\label{sec:phd-thesis-lineage-equation}

The Chatman Equation $A = \mu(O^*)$ formalizes the thesis that process authority is the
application of the process mining function to the full object-centric process state. It appears
in the thesis lineage recorded in \texttt{ledgers/PROJECT\_INDEX.yaml} and is operationalized
in the mathematical invariants of checkpoint \texttt{PROCESS\_INTELLIGENCE\_ALIVE\_001.md},
which lists the admissibility boundary as $R \vdash P_i = \mu(O^*, T, L)$, the autonomic
actuation function as $\alpha(K, P, L, T) \to \tau$, the token-game fitness formula
$f(L, N) = 1 - \frac{\sum m}{\sum c} - \frac{\sum r}{\sum p}$, and the OCPQ refinement
relation $p \sqsubseteq_L c$.

These formulas are not illustrative; they are the mathematical invariants verified in the
research program. The fitness formula is verified by the script
\texttt{experiments/validate\_log\_fitness.py}; the OCPQ refinement is verified by
\texttt{experiments/validate\_ocpq\_refinement.py}; the equations are verified by
\texttt{experiments/verify\_equations.py}. The Chatman Equation is thus grounded in executable
verification, not assertion alone.

The equation's significance is architectural: it specifies that authority over a process claim
can only be derived from a mining function applied to a specific, identified object-centric
process state --- not from a rule engine, not from a prediction model, and not from a
monitoring dashboard. This constraint is the source of the receipt doctrine and the named-witness
requirement that permeate the entire foundry.

\subsection{Relationship to Receipts and Replay}
\label{sec:phd-thesis-lineage-receipts}

Receipts and replay are the two mechanisms by which the Chatman Equation's authority constraint
is enforced at the artifact level. A receipt is a typed, witnessed, cryptographically bound
record that a specific process claim was derived from a specific event log using a specific
process model and algorithm. A replay is the execution of the token game against the event log
to compute fitness and precision metrics under a named algorithm.

The receipt doctrine is defined in \texttt{doctrine/RECEIPT\_DOCTRINE.md}. The physical receipt
chain is anchored by five JSON receipts in the \texttt{receipts/} directory:
\texttt{rec\_ebitda\_rework\_001.json} through \texttt{rec\_residual\_standard\_005.json},
each carrying an ed25519 validator signature, a target log hash, a process model hash, and
wasm4pm-computed verification results. For example, \texttt{rec\_ebitda\_rework\_001.json}
records: fitness 0.982, precision 0.945, throughput 4.12 days, EBITDA impact \$1,250,000 ---
all derived by the wasm4pm engine and bound to specific BLAKE3-intended hashes. The ed25519
signature is:
\texttt{ed25519\_sig\_8a3e811c\_02e0df...} (full value in the receipt file).

The ggen ecosystem manufactures these receipts through SPARQL queries over the process
intelligence ontology, rendering them via Tera templates into board-admissible claim surfaces.
The receipt chain is the physical implementation of the Chatman Equation's authority
requirement: every claim is $\mu(O^*, T, L)$ instantiated, hashed, signed, and registered.
